% Candidate rules: remote presentation (§4.2.1).  Written by hand; see docs/one-id-survey.md.

\section{Remote presentation}

\begin{definition}[The order of a verifier's steps over the DC API (C-45)]
  \label{def:c45}
  A verifier must do four things in order: detect the API, check which
  protocols the browser supports, request the credentials by calling the
  browser's credential request, and validate the response.

  A response is validated only after a request the same verifier made.

  External: W3C Digital Credentials API §2.1 to §2.3.

  Source:
  \href{https://standards.atlassian.net/wiki/spaces/OneID/pages/2947416082/Technical+Implementation+Guide}{One
  ID guide} §4.2.1.1.
\end{definition}

\begin{definition}[Supporting both protocols (C-46)]
  \label{def:c46}
  \uses{def:c45}
  A verifier should support both protocols, \texttt{org-iso-mdoc} and
  \texttt{openid4vp-v1-signed}, and it is recommended to send both in one
  request when the browser supports both.

  A request may carry two protocol requests. The wallet answers one of them.

  External: W3C Digital Credentials API §2.2, §5.

  Source:
  \href{https://standards.atlassian.net/wiki/spaces/OneID/pages/2947416082/Technical+Implementation+Guide}{One
  ID guide} §4.2.1.1.
\end{definition}

\begin{definition}[Registering the wallet with the operating system (C-48)]
  \label{def:c48}
  \notready
  A wallet provider must register its app with the operating system's
  credential manager. The prerequisite is marked informative and uses MUST.

  Outside the model.

  Open point: Is registration a requirement or informative?

  Source:
  \href{https://standards.atlassian.net/wiki/spaces/OneID/pages/2947416082/Technical+Implementation+Guide}{One
  ID guide} §4.2.1.2, prerequisite 1 and step 2.
\end{definition}

\begin{definition}[The reader's list of trusted issuers (C-49)]
  \label{def:c49}
  \uses{def:c15}
  \notready
  A reader must have access to lists of trusted issuer certificates. Without
  them validation fails.

  The issuer check takes the list as an input.

  Open point: The same requirement as C-15: which list, and how is it
  obtained?

  Source:
  \href{https://standards.atlassian.net/wiki/spaces/OneID/pages/2947416082/Technical+Implementation+Guide}{One
  ID guide} §4.2.1.2, prerequisite 2.
\end{definition}

\begin{definition}[Cipher suites (C-50)]
  \label{def:c50}
  Reader and wallet must conform to the cipher suites of ISO/IEC 18013-5.

  Outside the model: cryptography is taken as ideal.

  External: ISO/IEC 18013-5:2021 §9.1.5.2.

  Source:
  \href{https://standards.atlassian.net/wiki/spaces/OneID/pages/2947416082/Technical+Implementation+Guide}{One
  ID guide} §4.2.1.2, prerequisite 3.
\end{definition}

\begin{definition}[The wallet's check of a verifier's certificate (C-51)]
  \label{def:c51}
  \uses{def:c16}
  \notready
  The heading says a wallet must have a trusted list of verifiers. The text
  says it should have a way to check the authority that issued a verifier's
  certificate against a trusted list.

  Whether a wallet refuses an unknown verifier depends on C-16.

  Open point: Heading and text differ: must or should?

  Source:
  \href{https://standards.atlassian.net/wiki/spaces/OneID/pages/2947416082/Technical+Implementation+Guide}{One
  ID guide} §4.2.1.2, prerequisite 4.
\end{definition}

\begin{definition}[Several document types in one request (C-53)]
  \label{def:c53}
  \uses{def:mdoc-request, def:photoid, def:boarding-pass}
  One request may name several document types, such as the photoid and the
  boarding pass.

  This is what 'presented together' means in the model: one response to one
  request that names both.

  Source:
  \href{https://standards.atlassian.net/wiki/spaces/OneID/pages/2947416082/Technical+Implementation+Guide}{One
  ID guide} §4.2.1.2, step 1.
\end{definition}

\begin{definition}[Requesting only what is needed (C-54)]
  \label{def:c54}
  \uses{def:c11, def:mdoc-request}
  \notready
  A relying party should request only the namespaces and data elements it
  needs.

  Same content as C-11 at a weaker level.

  Open point: §3.1 says must, this step says should. See C-11.

  Source:
  \href{https://standards.atlassian.net/wiki/spaces/OneID/pages/2947416082/Technical+Implementation+Guide}{One
  ID guide} §4.2.1.2, step 1.
\end{definition}

\begin{definition}[Reader authentication in the mdoc request (C-55)]
  \label{def:c55}
  \uses{def:c16, def:mdoc-request}
  \notready
  Reader authentication is recommended so the wallet can tell the request
  comes from a valid verifier. The guide notes not every wallet can process
  it, since it is optional in the specification, and says a verifier should
  take that into account.

  If a request need not be authenticated, a wallet may answer any verifier,
  and consent is the only gate.

  Open point: See C-16: is a request without reader authentication answered?

  External: ISO/IEC TS 23220-4:2026 Annex A §A.1.4.

  Source:
  \href{https://standards.atlassian.net/wiki/spaces/OneID/pages/2947416082/Technical+Implementation+Guide}{One
  ID guide} §4.2.1.2, step 1.
\end{definition}

\begin{definition}[The wallet's check of the request's origin (C-57)]
  \label{def:c57}
  \uses{def:c16, def:c55}
  \notready
  A wallet should check that a request comes from a trusted verifier, by
  validating reader authentication and the session transcript of the request.

  With 'should', a wallet that skips the check still conforms.

  Open point: See C-16.

  External: ISO/IEC 18013-5:2021 §9.1.4, §9.1.5.

  Source:
  \href{https://standards.atlassian.net/wiki/spaces/OneID/pages/2947416082/Technical+Implementation+Guide}{One
  ID guide} §4.2.1.2, step 3.
\end{definition}

\begin{definition}[The wallet's check of the request's content (C-58)]
  \label{def:c58}
  \uses{def:mdoc-request}
  A wallet should validate the request's syntax, document type, namespaces
  and data element identifiers, and may apply its own disclosure policies to
  accept or refuse the request.

  A wallet answers only for data elements it holds.

  Source:
  \href{https://standards.atlassian.net/wiki/spaces/OneID/pages/2947416082/Technical+Implementation+Guide}{One
  ID guide} §4.2.1.2, step 4.
\end{definition}

\begin{definition}[The mdoc response (C-59)]
  \label{def:c59}
  \uses{def:mdoc-request}
  The response must follow Annex C: it names the protocol \texttt{dcapi}, is
  encrypted to the verifier's key, and uses the session transcript of Annex
  C.

  A response is readable only by the verifier whose key the request carried,
  and its transcript ties it to that request.

  External: ISO/IEC TS 18013-7:2025 Annex C §C.3, §C.4, §C.5; IETF RFC 9180.

  Source:
  \href{https://standards.atlassian.net/wiki/spaces/OneID/pages/2947416082/Technical+Implementation+Guide}{One
  ID guide} §4.2.1.2, step 4.
\end{definition}

\begin{definition}[Decrypting the mdoc response (C-61)]
  \label{def:c61}
  \uses{def:c59}
  The verifier must decrypt the response with the mechanism of Annex C.

  Only the verifier that made the request can read the response.

  External: ISO/IEC TS 18013-7:2025 Annex C §C.4; IETF RFC 9180.

  Source:
  \href{https://standards.atlassian.net/wiki/spaces/OneID/pages/2947416082/Technical+Implementation+Guide}{One
  ID guide} §4.2.1.2, step 6.
\end{definition}

\begin{definition}[The passport security object check (C-64)]
  \label{def:c64}
  \uses{thm:portrait-is-dg2, def:photoid-advisories}
  \notready
  When a photoid carries data groups and the passport's security object, a
  verifier may check the security object as ICAO Doc 9303 Part 11 describes.

  The guide says the data groups namespace binds the credential to the
  passport chip's content. The check covers the data groups; whether it
  covers the portrait a touchpoint matches depends on the portrait being the
  same as data group 2.

  Open point: When is the check expected, and does it cover the portrait a
  touchpoint matches against?

  External: ICAO Doc 9303 Part 11 §5.

  Source:
  \href{https://standards.atlassian.net/wiki/spaces/OneID/pages/2947416082/Technical+Implementation+Guide}{One
  ID guide} §4.2.1.2 step 6, §4.2.1.3 step 5.
\end{definition}

\begin{definition}[What the mdoc profile does not check (C-65)]
  \label{def:c65}
  \uses{def:mdoc-request}
  \notready
  The verifier's list of checks in the mdoc profile has no step that compares
  the response with the request: that it holds only the document types and
  data elements requested, and that it answers this request's nonce. The
  OpenID4VP profile states a nonce check and a check that the token matches
  the query, and the proximity profile says issuer data validation covers the
  data elements requested.

  If the comparison is only implied by the ISO clauses, the model must assume
  it. If it is not, a response with more data than requested, or made for
  another request, is accepted.

  Case to check: A relying party requests the family name and receives a
  valid response that also holds the date of birth.

  Open point: Must a verifier in the mdoc profile check that the response
  matches its request and its nonce, as the OpenID4VP profile requires?

  Lookup: ISO/IEC TS 18013-7:2025 Annex C §C.3 to §C.5; ISO/IEC 18013-5:2021
  §9.1.3.

  Source:
  \href{https://standards.atlassian.net/wiki/spaces/OneID/pages/2947416082/Technical+Implementation+Guide}{One
  ID guide} §4.2.1.2, step 6.
\end{definition}

\begin{definition}[Reporting the result (C-66)]
  \label{def:c66}
  \uses{def:acceptance}
  \notready
  After processing a response, 'successful or error handled', the verifier
  must set the error structure that the ISO device response defines, and
  should return that standard's error codes.

  Not modelled.

  Open point: The structure cited is part of the wallet's response. What does
  a verifier set, and to whom does it return codes?

  External: ISO/IEC 18013-5:2021 §8.3.2.1.2.2, §8.3.2.1.2.3.

  Source:
  \href{https://standards.atlassian.net/wiki/spaces/OneID/pages/2947416082/Technical+Implementation+Guide}{One
  ID guide} §4.2.1.2 step 7, §4.2.2 step 6.
\end{definition}

\begin{definition}[Passing validated data on (C-67)]
  \label{def:c67}
  \uses{def:acceptance}
  After a successful validation the verifier should send the data to the
  identity management system through IATA's CUWS APIs.

  Data leaves the verifier only after acceptance.

  External: IATA CUWS standard APIs.

  Source:
  \href{https://standards.atlassian.net/wiki/spaces/OneID/pages/2947416082/Technical+Implementation+Guide}{One
  ID guide} §4.2.1.2 step 7, §4.2.1.3 step 6, §4.2.2 step 6.
\end{definition}

\begin{definition}[Complying with the high assurance profile (C-68)]
  \label{def:c68}
  \notready
  Verifier and wallet must comply with the OpenID4VC high assurance profile
  1.0 and follow Appendix A of OpenID4VP 1.0, and the profile's security and
  privacy considerations take precedence over the appendix's. Step 2 then
  says the request should comply with the profile's sections.

  Taken as must. For the DC API the profile's own considerations point to
  those of the appendix (HAIP §10, §11), so the precedence the guide states
  changes nothing.

  Open point: Must or should: the opening paragraph and step 2 differ.

  External: OpenID4VC HAIP 1.0 §5, §5.2, §5.3.1, §10, §11; OpenID4VP 1.0
  Appendix A (read).

  Source:
  \href{https://standards.atlassian.net/wiki/spaces/OneID/pages/2947416082/Technical+Implementation+Guide}{One
  ID guide} §4.2.1.3.
\end{definition}

\begin{definition}[The credential query (C-69)]
  \label{def:c69}
  \uses{def:element, def:c11}
  \notready
  The verifier must write a DCQL query: the credentials, each with the format
  \texttt{mso\_mdoc} and a document type, and the claims wanted, each a path
  of a namespace then a data element name. The claims must be the minimum
  needed. A claim may say whether the relying party intends to retain it;
  that member is optional.

  A query names data elements exactly as the mdoc request does, so one notion
  of request serves both profiles. In the guide's own first example one
  namespace is written with a leading space, which names no namespace of a
  photoid.

  Open point: What is the minimum set of claims? The same question as C-11,
  here stated with must.

  External: OpenID4VP 1.0 §6, §B.2.4.

  Source:
  \href{https://standards.atlassian.net/wiki/spaces/OneID/pages/2947416082/Technical+Implementation+Guide}{One
  ID guide} §4.2.1.3, step 1.
\end{definition}

\begin{definition}[Trusted authorities in the query (C-71)]
  \label{def:c71}
  \uses{def:c69, def:c15}
  Verifier and wallet must support the trusted authorities query, with which
  a verifier names the authorities whose credentials it accepts.

  A wallet can pick a credential the verifier will trust.

  External: OpenID4VP 1.0 §6.1.1.1; OpenID4VC HAIP 1.0 §5.

  Source:
  \href{https://standards.atlassian.net/wiki/spaces/OneID/pages/2947416082/Technical+Implementation+Guide}{One
  ID guide} §4.2.1.3, step 2.
\end{definition}

\begin{definition}[The nonce and state of the authorization request (C-72)]
  \label{def:c72}
  \notready
  The nonce and state must be set as OpenID4VP says. OpenID4VP requires a
  nonce that is fresh and random for every request, and asks for 'sufficient
  entropy' without a number (OpenID4VP §5.2). It does not define
  \texttt{state} over the DC API, where a wallet ignores it (OpenID4VP §5.3,
  §A.2). The guide's example request has a seven-digit nonce and no state,
  and the mdoc profile requires 16 bytes of entropy.

  A nonce is made fresh for every request, so a response made for one request
  does not answer another, provided the nonce cannot be guessed.

  Open point: What is the minimum entropy of the nonce in this profile: the 16
  bytes the mdoc profile requires? And is the must for \texttt{state} meant,
  given that it has no effect over the DC API?

  External: OpenID4VP 1.0 §5.2, §5.3, §A.2 (read).

  Source:
  \href{https://standards.atlassian.net/wiki/spaces/OneID/pages/2947416082/Technical+Implementation+Guide}{One
  ID guide} §4.2.1.3, step 2.
\end{definition}

\begin{definition}[Consent in the OpenID4VP profile (C-75)]
  \label{def:c75}
  \uses{def:c22, def:openid4vp-wallet-checks, def:openid4vp-binding}
  \notready
  Informed consent must be taken and the wallet must show what data is shared
  with whom. After user authentication and authorization the response
  'should' be released.

  Taken with §3.2: no response before authentication and consent.

  Open point: §3.2 says a presentation must happen only after authentication
  and consent. Does 'should be released' weaken that, or only describe the
  release?

  Source:
  \href{https://standards.atlassian.net/wiki/spaces/OneID/pages/2947416082/Technical+Implementation+Guide}{One
  ID guide} §4.2.1.3, step 4.
\end{definition}

\begin{definition}[Issuer and device checks in the OpenID4VP profile (C-78)]
  \label{def:c78}
  \uses{def:c15, def:c03, def:expired, def:c101}
  \notready
  The verifier must authenticate the issuer against a trusted list, which it
  'should' have access to, authenticate the issuer data, and authenticate the
  mdoc, which the guide calls holder binding. This profile's text does not
  name a time validity check (C-101).

  An accepted response comes from a trusted issuer and the holder's device.
  That it is not expired depends on C-101.

  Open point: §3.2 says a relying party must have the list of trusted
  issuers, this step says should. Which holds?

  External: ISO/IEC 18013-5:2021 §9.1.2, §9.1.3; ISO/IEC TS 23220-4:2026
  Annex A §A.3.1.

  Source:
  \href{https://standards.atlassian.net/wiki/spaces/OneID/pages/2947416082/Technical+Implementation+Guide}{One
  ID guide} §4.2.1.3, step 5.
\end{definition}

\begin{definition}[Errors in the OpenID4VP profile (C-79)]
  \label{def:c79}
  \uses{def:acceptance}
  Errors in processing a response should follow the OpenID4VP error
  responses.

  Not modelled.

  External: OpenID4VP 1.0 §8.5.

  Source:
  \href{https://standards.atlassian.net/wiki/spaces/OneID/pages/2947416082/Technical+Implementation+Guide}{One
  ID guide} §4.2.1.3, step 6.
\end{definition}

\begin{definition}[When the browser lacks the API or the protocol (C-99)]
  \label{def:c99}
  \uses{def:c45}
  \notready
  The verifier must check that the browser supports the API and the
  protocols. The guide does not say what the verifier does when it does not.

  A remote presentation has no defined outcome on a browser without support.

  Open point: What does a relying party offer a passenger whose browser or
  wallet supports neither protocol?

  Source:
  \href{https://standards.atlassian.net/wiki/spaces/OneID/pages/2947416082/Technical+Implementation+Guide}{One
  ID guide} §4.2.1.1.
\end{definition}

\begin{definition}[The version of the Digital Credentials API (C-100)]
  \label{def:c100}
  \uses{def:c45}
  \notready
  The guide puts a dated draft of the W3C Digital Credentials API in scope
  and notes the standard is still a draft.

  The protocol identifiers and the API are taken from that draft. A later
  draft may change them.

  Open point: When the W3C draft changes, which version must verifiers and
  wallets support, and who decides?

  External: W3C Digital Credentials API, draft of 16 July 2026.

  Source:
  \href{https://standards.atlassian.net/wiki/spaces/OneID/pages/2947416082/Technical+Implementation+Guide}{One
  ID guide} §1.4, §4.2.1.1.
\end{definition}

\begin{definition}[Whether issuer data authentication includes the validity check (C-101)]
  \label{def:c101}
  \uses{def:expired}
  \notready
  The mdoc profile cites one clause of ISO/IEC 18013-5 for issuer data
  authentication and says it covers time validity. The OpenID4VP profile and
  the proximity profile cite a different clause with must, and do not say
  whether it covers time validity; proximity adds the first clause with
  should.

  If the clause cited with must includes the validity check, an expired
  credential is refused in all three profiles. If it does not, two profiles
  accept one.

  Open point: Does the clause the OpenID4VP and proximity profiles cite with
  must include the time validity check?

  Lookup: ISO/IEC 18013-5:2021 §9.1.2 and §9.3.1.

  External: ISO/IEC 18013-5:2021 §9.1.2, §9.3.1.

  Source:
  \href{https://standards.atlassian.net/wiki/spaces/OneID/pages/2947416082/Technical+Implementation+Guide}{One
  ID guide} §4.2.1.3 step 5, §4.2.2 step 5.
\end{definition}

\begin{definition}[Checking the data model on what is received (C-107)]
  \label{def:c107}
  \uses{def:acceptance, def:photoid, def:boarding-pass}
  \notready
  The guide's lists of a relying party's checks hold no step that checks a
  response against the credential data model of §4.1: the mandatory data
  elements, their encodings, the format of the boarding pass string.  With
  selective disclosure a relying party sees only the data elements it
  requested.

  A relying party cannot tell from a response whether the credential behind
  it holds every mandatory data element.  The data model requirements can be
  checked in full only by the issuer, and by a wallet that holds the whole
  credential.

  Case to check: A relying party requests the family name of a photoid.  The
  photoid behind the response has no portrait, which one reading of the guide
  requires.  The response is the same as one from a complete photoid.

  Open point: Does a relying party check anything of §4.1 on a response, such
  as the encoding of what it received or the format of the boarding pass
  string?  Or are the data model requirements the issuer's duty alone?

  Source:
  \href{https://standards.atlassian.net/wiki/spaces/OneID/pages/2947416082/Technical+Implementation+Guide}{One ID guide} §4.1.1, §4.1.2, §4.2.1.3.
\end{definition}

% The presentation rules in Lean.

\begin{definition}[What a wallet and a relying party exchange]
  \label{def:exchange}
  \lean{Rules.OneId.Request, Rules.OneId.Consent, Rules.OneId.Response, Rules.OneId.Presented, Rules.OneId.Credential, Rules.OneId.ElementRef, Rules.OneId.Profile, Rules.OneId.SessionTranscript, Rules.OneId.DeviceSignature, Rules.OneId.Key, Rules.OneId.Nonce, Rules.OneId.Origin, Rules.OneId.Issuer, Rules.OneId.Mdoc.refs, Rules.OneId.Mdoc.disclose, Rules.OneId.Response.refs}
  \uses{def:mdoc, def:element}
  All three profiles have the same parts: a request that names data elements by
  document type, namespace and identifier, the holder's consent, a response
  that discloses data elements of credentials, and the relying party's checks.
  A credential as issued has an issuer and the key of the device it was issued
  to.  A presented credential has the data elements disclosed, the issuer's
  signature and a device signature.

  No cryptography is modelled.  Keys, nonces and origins are values that can
  only be compared, and a device signature is the key that made it and the
  session transcript it was made over.

  Source: \href{https://standards.atlassian.net/wiki/spaces/OneID/pages/2947416082/Technical+Implementation+Guide}{One ID guide} §4.2.1.1, §4.2.1.2, §4.2.1.3, §4.2.2.
\end{definition}

\begin{definition}[The session transcript of each profile]
  \label{def:session-transcript}
  \lean{Rules.OneId.IsoTranscripts, Rules.OneId.expectedTranscript}
  \uses{def:exchange, def:openid4vp-binding, def:c59}
  A device signs its response over a session transcript, and each profile
  defines the transcript differently.  The OpenID4VP profile's is in a public
  specification.  The other two are in ISO clauses this project has not read,
  so the model takes them as given: for a request, some value, with no
  assumption about which parts of the request it depends on.

  External: ISO/IEC TS 18013-7:2025 Annex C §C.5; ISO/IEC 18013-5:2021
  §9.1.5.1.

  Source: \href{https://standards.atlassian.net/wiki/spaces/OneID/pages/2947416082/Technical+Implementation+Guide}{One ID guide} §4.2.1.2, §4.2.2.
\end{definition}

\begin{definition}[A wallet's answer to a request (C-60)]
  \label{def:wallet-respond}
  \lean{Rules.OneId.Wallet.respond, Rules.OneId.Credential.present, Rules.OneId.ConsentReading, Rules.OneId.Request.wanted}
  \uses{def:exchange, def:session-transcript, def:c22}
  A wallet answers a request with the data elements that the request names and
  the holder approved, and with nothing of a credential that has none of them.
  It signs each credential it presents with that credential's device key, over
  the session transcript of the request as the wallet received it.

  The guide requires informed consent before every response, in all three
  profiles.  What a wallet does when the holder approves only some of the data
  elements requested is a reading: it sends those, or it sends nothing of that
  credential.

  External: OpenID4VP §6.4; ISO/IEC 18013-5:2021 §8.3.2.1.2.2, §9.1.3.

  Source: \href{https://standards.atlassian.net/wiki/spaces/OneID/pages/2947416082/Technical+Implementation+Guide}{One ID guide} §3.1, §3.2, §4.2.1.2, §4.2.1.3, §4.2.2.
\end{definition}

\begin{theorem}[A wallet discloses only what was requested and approved]
  \label{thm:wallet-discloses}
  \lean{Rules.OneId.wallet_discloses_only_requested_and_approved, Rules.OneId.wallet_presents_nothing_without_consent, Rules.OneId.wallet_signs_with_device_key}
  \uses{def:wallet-respond}
  Every data element a wallet discloses was named by the request and approved
  by the holder, under both readings of partial approval.  A wallet presents
  nothing when the holder approved nothing, and signs each credential it
  presents with that credential's device key over the transcript of the request
  it received.
\end{theorem}
\begin{proof}
  \uses{def:wallet-respond}
  The wallet discloses the requested data elements filtered by the approved ones.
\end{proof}

\begin{theorem}[Partial approval withholds the credential (C-105)]
  \label{thm:partial-consent}
  \lean{Rules.OneId.partial_consent_withholds_credential}
  \uses{def:wallet-respond}
  A wallet presents nothing of a credential when the holder approved only some
  of the data elements requested from it.  Not proven.

  Counterexample: a check-in page requests the family name and the portrait.
  The traveler approves the family name and declines the portrait.  A wallet
  that sends what was approved sends the family name.

  Open point: the guide requires informed consent and that the wallet show what
  is shared with whom, and does not say what happens when a holder agrees to a
  part.  OpenID4VP says a wallet that cannot deliver every claim asked of a
  credential must not return it (OpenID4VP §6.4.1).  When a holder approves
  only some of the data elements requested, does the wallet send those, or
  nothing from that credential?  Is the answer the same in all three profiles?

  Source: \href{https://standards.atlassian.net/wiki/spaces/OneID/pages/2947416082/Technical+Implementation+Guide}{One ID guide} §3.2, §4.2.1.3.
\end{theorem}

\begin{definition}[A relying party's decision on a response (C-47)]
  \label{def:acceptance}
  \lean{Rules.OneId.refusals, Rules.OneId.Presented.refusals, Rules.OneId.AcceptanceReading, Rules.OneId.AcceptanceReading.checksValidity, Rules.OneId.AcceptanceReading.checksTranscript, Rules.OneId.Refusal}
  \uses{def:exchange, def:session-transcript, def:expired, def:c03, def:c13, def:c15, def:c45, def:c49, def:c61, def:c65, def:c78, def:c88, def:c101}
  A relying party accepts a response only if every credential in it comes from
  an issuer on its list of trusted issuers, has an issuer's signature that
  verifies, and has a device signature made with the device key the issuer
  signed.  All three profiles require these.

  The profiles differ on three more checks, and the differences are readings:
  whether the device signature must be over the transcript of this request
  (required in the OpenID4VP profile), whether the end of validity is checked
  (required in the mdoc profile), and whether a response with a data element
  that was not requested is refused (stated nowhere).  Who is on the list of
  trusted issuers is given to the rule.  Decryption and revocation are not
  modelled.

  External: ISO/IEC 18013-5:2021 §9.1.2, §9.1.3, §9.3.1; ISO/IEC TS
  18013-7:2025 Annex C; OpenID4VP §8.6.

  Source: \href{https://standards.atlassian.net/wiki/spaces/OneID/pages/2947416082/Technical+Implementation+Guide}{One ID guide} §3.1, §3.2, §4.2.1.1, §4.2.1.2, §4.2.1.3, §4.2.2.
\end{definition}

\begin{theorem}[An accepted response comes from a trusted issuer]
  \label{thm:accepted-trusted-issuer}
  \lean{Rules.OneId.accepted_from_trusted_issuer}
  \uses{def:acceptance}
  Every credential in an accepted response comes from an issuer on the relying
  party's list and has an issuer's signature that verifies, in every profile
  and under every reading.
\end{theorem}
\begin{proof}
  \uses{def:acceptance}
  An issuer off the list and a signature that does not verify are each a reason to refuse.
\end{proof}

\begin{theorem}[An accepted response is signed with the device key]
  \label{thm:accepted-device-key}
  \lean{Rules.OneId.accepted_signed_by_device_key}
  \uses{def:acceptance}
  Every credential in an accepted response has a device signature made with the
  device key its issuer signed, in every profile and under every reading.  This
  is the device half of the guarantee G-02.
\end{theorem}
\begin{proof}
  \uses{def:acceptance}
  A missing device signature and one made with another key are each a reason to refuse.
\end{proof}

\begin{theorem}[A wallet cannot be checked for consent]
  \label{thm:no-check-shows-consent}
  \lean{Rules.OneId.no_check_shows_consent}
  \uses{def:acceptance, def:wallet-respond}
  Whatever a relying party checks on a request and its response, it cannot tell
  whether the holder approved what was sent.  The consent half of the guarantee
  G-01 rests on the wallet, and on whoever vouches for wallets.
\end{theorem}
\begin{proof}
  \uses{def:wallet-respond}
  The same response could have been sent with nothing approved.
\end{proof}

\begin{theorem}[An honest response is accepted]
  \label{thm:wallet-response-accepted}
  \lean{Rules.OneId.wallet_response_accepted}
  \uses{def:acceptance, def:wallet-respond}
  A response a wallet makes for the relying party's own request, from
  credentials of trusted issuers that have not expired, is accepted in every
  profile and under every reading.
\end{theorem}
\begin{proof}
  \uses{def:acceptance, def:wallet-respond}
  Each reason to refuse fails for such a response.
\end{proof}

\begin{definition}[A photoid and a boarding pass presented together (C-106)]
  \label{def:pair}
  \lean{Rules.OneId.pairRefusals, Rules.OneId.pairMatch, Rules.OneId.PairReading, Rules.OneId.PairRefusal, Rules.OneId.Presented.isOf}
  \uses{def:acceptance, def:traveler-match, def:c03, def:c53}
  A relying party accepts a photoid and a boarding pass presented together only
  if each is accepted on its own and the two pass a check that they are one
  traveler's.  The guide states no such check.  The content of the two can be
  compared, as in the credential data model, and a presentation adds the device
  keys: a reading in which the two credentials must carry the same device key.

  External: IATA Resolution 792; ISO/IEC 18013-5:2021 §9.1.3.

  Source: \href{https://standards.atlassian.net/wiki/spaces/OneID/pages/2947416082/Technical+Implementation+Guide}{One ID guide} §4.1.2, §3.2, §4.2.1.2.
\end{definition}

\begin{theorem}[In an accepted pair each credential is accepted]
  \label{thm:pair-each-accepted}
  \lean{Rules.OneId.accepted_pair_each_accepted}
  \uses{def:pair}
  In an accepted pair each credential is accepted on its own, the first is a
  photoid and the second is a boarding pass.
\end{theorem}
\begin{proof}
  \uses{def:pair}
  Each reason to refuse one credential is a reason to refuse the pair.
\end{proof}

\begin{theorem}[No check of a presentation identifies the traveler]
  \label{thm:no-presentation-check}
  \lean{Rules.OneId.no_presentation_check_identifies_the_traveler}
  \uses{def:pair, thm:no-content-check}
  Whatever a relying party checks on a presented photoid and a presented
  boarding pass, device keys included, a pair it accepts can be presentations
  of credentials issued about two different travelers.  A check of device keys
  gives one traveler only together with a rule about issuance.
\end{theorem}
\begin{proof}
  \uses{def:pair}
  The same two presentations could come from credentials issued about different people.
\end{proof}

\begin{theorem}[One traveler's pair is accepted]
  \label{thm:same-traveler-pair-accepted}
  \lean{Rules.OneId.same_traveler_pair_accepted}
  \uses{def:pair, thm:same-traveler-matches}
  A photoid and a boarding pass about the same traveler, each accepted on its
  own, are accepted together.  Not proven.

  Counterexample: a traveler holds a photoid on a phone and a boarding pass on
  a tablet.  Under the reading that the two credentials must carry the same
  device key, the pair is refused.

  Open point: if the device key is what ties the two credentials together, a
  traveler must hold both in one wallet on one device, and both issuers must
  issue to the same key.  Does One ID require a photoid and a boarding pass to
  be issued to the same device key, and what does a traveler do who holds them
  on two devices?

  Source: \href{https://standards.atlassian.net/wiki/spaces/OneID/pages/2947416082/Technical+Implementation+Guide}{One ID guide} §4.1.2, §3.2.
\end{theorem}

\begin{definition}[The OpenID4VP request (C-70)]
  \label{def:openid4vp-request}
  \lean{Rules.OneId.OpenId4Vp.requestFindings, Rules.OneId.OpenId4Vp.AuthorizationRequest, Rules.OneId.OpenId4Vp.CredentialQuery, Rules.OneId.OpenId4Vp.Certificate, Rules.OneId.OpenId4Vp.RequestReading, Rules.OneId.OpenId4Vp.RequestFinding, Rules.OneId.OpenId4Vp.AuthorizationRequest.shared, Rules.OneId.OpenId4Vp.protocolValue, Rules.OneId.OpenId4Vp.responseTypeValue, Rules.OneId.OpenId4Vp.responseModeValue, Rules.OneId.OpenId4Vp.clientIdPrefixValue, Rules.OneId.OpenId4Vp.mdocFormat}
  \uses{def:exchange, def:c45, def:c68, def:c69, def:c72}
  An OpenID4VP request is accepted only if it has the protocol value, response
  type, response mode and client identifier prefix the guide names, is signed
  with the key of the certificate its client identifier names, names its
  expected origins, carries an encryption key and a nonce, and asks only for
  mdoc credentials.  Whether it must also carry a state value is a reading.

  The rule does not check that the key and the nonce are new for each request,
  the minimum set of claims (C-69), or the entropy of the nonce (C-72).

  External: OpenID4VP §5, §5.1, §5.2, §5.9.3, §6, §A.2, §A.3.2; OpenID4VC HAIP
  §5, §5.2; W3C Digital Credentials API §5.

  Source: \href{https://standards.atlassian.net/wiki/spaces/OneID/pages/2947416082/Technical+Implementation+Guide}{One ID guide} §4.2.1.3.
\end{definition}

\begin{theorem}[An accepted OpenID4VP request has its required parts]
  \label{thm:openid4vp-request-parts}
  \lean{Rules.OneId.OpenId4Vp.accepted_request_parts}
  \uses{def:openid4vp-request}
  An accepted request has the four values the guide names, is signed with the
  key of the certificate its client identifier names, names an expected origin,
  carries an encryption key and a nonce, and asks only for mdoc credentials.
\end{theorem}
\begin{proof}
  \uses{def:openid4vp-request}
  Each missing part is a finding.
\end{proof}

\begin{theorem}[An accepted OpenID4VP request carries a state value]
  \label{thm:openid4vp-state}
  \lean{Rules.OneId.OpenId4Vp.accepted_request_has_state}
  \uses{def:openid4vp-request, def:c72}
  An accepted OpenID4VP request carries a state value.  Not proven.

  Counterexample: a signed request with every part the guide's step 2 lists and
  no state value, as in the guide's own example, is accepted under the reading
  that follows OpenID4VP.

  Open point: the guide says the nonce and state must be set as OpenID4VP §5.2
  says.  That clause requires state only for a presentation without a holder
  binding proof, §5.3 exempts the Digital Credentials API, and Appendix A, which
  the guide says must be followed, does not list the parameter.  Is the must
  for state meant, and if so what does a relying party or a wallet do with it?

  Source: \href{https://standards.atlassian.net/wiki/spaces/OneID/pages/2947416082/Technical+Implementation+Guide}{One ID guide} §4.2.1.3.
\end{theorem}

\begin{definition}[The wallet's checks of an OpenID4VP request (C-73)]
  \label{def:openid4vp-wallet-checks}
  \lean{Rules.OneId.OpenId4Vp.walletFindings, Rules.OneId.OpenId4Vp.WalletFinding}
  \uses{def:openid4vp-request, def:c16}
  A wallet answers an OpenID4VP request only if the request is well formed,
  came from an origin it names, is signed with the key of the certificate its
  client identifier names, and that certificate is a trusted relying party's.

  The guide introduces its list of four checks with should and marks three of
  them with must.  OpenID4VP, which the guide makes mandatory, requires the
  origin check and the check of the signature and the certificate chain of a
  wallet, so the model treats the checks as required.  Which list a wallet
  checks a relying party against is given to the rule (C-16).

  External: OpenID4VP §A.2, §5.9.3.

  Source: \href{https://standards.atlassian.net/wiki/spaces/OneID/pages/2947416082/Technical+Implementation+Guide}{One ID guide} §4.2.1.3, §3.2.
\end{definition}

\begin{theorem}[A wallet answers only a trusted, signed request from an expected origin]
  \label{thm:openid4vp-answered}
  \lean{Rules.OneId.OpenId4Vp.answered_request_parts}
  \uses{def:openid4vp-wallet-checks}
  A request a wallet answers came from an origin it names, is signed with the
  key of the certificate its client identifier names, and that certificate is
  on the wallet's list of trusted relying parties.
\end{theorem}
\begin{proof}
  \uses{def:openid4vp-wallet-checks}
  Each failed check is a finding.
\end{proof}

\begin{definition}[What an OpenID4VP response is bound to (C-77)]
  \label{def:openid4vp-binding}
  \lean{Rules.OneId.OpenId4Vp.transcript, Rules.OneId.OpenId4Vp.deviceResponses}
  \uses{def:exchange}
  In the OpenID4VP profile the wallet signs its response over the origin of the
  request, the nonce and the relying party's encryption key, and the relying
  party rebuilds that transcript and validates the nonce.  Each mdoc is
  returned in a device response of its own.

  The specification hashes the three values.  The model keeps the values, so it
  does not model two inputs with one hash.

  External: OpenID4VP §B.2.5, §B.2.6.2; OpenID4VC HAIP §5.3.1.

  Source: \href{https://standards.atlassian.net/wiki/spaces/OneID/pages/2947416082/Technical+Implementation+Guide}{One ID guide} §4.2.1.3.
\end{definition}

\begin{theorem}[An OpenID4VP response cannot be used for another request]
  \label{thm:openid4vp-refuses-other-request}
  \lean{Rules.OneId.OpenId4Vp.refuses_other_request, Rules.OneId.OpenId4Vp.accepted_for_one_request, Rules.OneId.OpenId4Vp.one_mdoc_per_deviceResponse}
  \uses{def:openid4vp-binding, def:acceptance}
  In the OpenID4VP profile a response signed for another nonce, another origin
  or another encryption key is refused, under every reading.  A response
  accepted for two requests was made for one nonce, so a captured response
  cannot answer a later request.
\end{theorem}
\begin{proof}
  \uses{def:openid4vp-binding, def:acceptance}
  The transcript the relying party rebuilds differs when any of the three values differs.
\end{proof}

\begin{definition}[The request of the mdoc profile (C-52)]
  \label{def:mdoc-request}
  \lean{Rules.OneId.MdocProfile.requestFindings, Rules.OneId.MdocProfile.DeviceRequest, Rules.OneId.MdocProfile.RequestFinding, Rules.OneId.MdocProfile.DeviceRequest.shared, Rules.OneId.MdocProfile.protocolValue, Rules.OneId.MdocProfile.encryptionProtocolValue, Rules.OneId.MdocProfile.minimumNonceEntropy}
  \uses{def:exchange, def:c45, def:element}
  A request of the mdoc profile is accepted only if it has the protocol value
  \texttt{org-iso-mdoc}, its encryption information names \texttt{dcapi}, and it
  carries a nonce of at least 16 bytes of entropy and the recipient's public
  key.  The entropy of a nonce is taken as given.  Reader authentication is
  recommended and is not checked by this rule.

  External: ISO/IEC TS 18013-7:2025 Annex C §C.2; W3C Digital Credentials API
  §5.

  Source: \href{https://standards.atlassian.net/wiki/spaces/OneID/pages/2947416082/Technical+Implementation+Guide}{One ID guide} §4.2.1.2.
\end{definition}

\begin{theorem}[An accepted mdoc request has its required parts]
  \label{thm:mdoc-request-parts}
  \lean{Rules.OneId.MdocProfile.accepted_deviceRequest_parts}
  \uses{def:mdoc-request}
  An accepted request of the mdoc profile names both protocol values and
  carries a nonce of at least 16 bytes of entropy and a recipient key.
\end{theorem}
\begin{proof}
  \uses{def:mdoc-request}
  Each missing part is a finding.
\end{proof}
